%%%
%
% $Autor: Wings $
% $Datum: 2021-05-14 $
% $Pfad: GitLab/MLEdgeComputer $
% $Dateiname: PowerSource
% $Version: 4620 $
%
% !TeX spellcheck = de_DE/GB
% !TeX program = pdflatex
% !BIB program = biber/bibtex
% !TeX encoding = utf8
%
%%%



\chapter{Power Source}

For some application, using the computer to power the device - although most common - would not work. For better results in recording subjects' everyday moves this device is designed; therefore, it requires a light-weighted, long-term portability powering system, and is easily carried in a pocket.

Commercially available with various power densities and voltages, batteries have gained popularity \cite{Dewan:2014}. Considering various aspects can make choosing the right type of battery a challenging task.

Understanding your device's power needs is essential to selecting an appropriate battery. To fulfill their requirements for higher outputs of electricity, devices such as EVs must rely on batteries of greater energy density and more powerful outputs.

The dimensions and mass of the battery hinge on those of the device it powers. Smartphones and laptops require batteries that are both lightweight and compact.

How often the device is used will dictate the necessary battery life. The duration between charges will also play a role. Devices that require frequent use should have batteries with extended cycle lives.

Factors like temperature and humidity, which fall under environmental conditions, may impact how effectively the battery operates. Special considerations should be made when planning on using batteries in areas that have extreme temperatures or humid conditions.

Selecting an appropriate battery requires consideration of potential safety risks. Certain types like lithium-ion have known concerns and need to be taken into account. Prioritizing safety concerns when making decisions is important.

The power supply for the device should be lightweight, portable, and long-lasting.
The power supply should also not add excessive weight to the device. So the lighter the power supply is the better. \cite{Scherer:2022}

\section{Port}

All the Arduino boards need power to operate, either it comes from the USB connection with Laptop, Ac power adapter, Battery or a regulated power supply. 

The power supplies for the Arduino Nano 33 BLE Sense board are connected via:

\begin{itemize}
    \item MicroUSB port
    \item Pins: Vin, GND
\end{itemize}

As a result, there are two ways to power up the Arduino Nano 33 BLE Sense board. The first option is to use the Pins: Pin 15 (Vin), with an input voltage range of 6.2V to 20V (optimal range of 7V to 12V) for the positive (+) connection and Pin 14 as GND for the negative connection (-) (See \ref{Schnittstellen} for the PIN Configuration). This option allows the USB port to remain free.


The second option is to use the MicroUSB port with either a lithium battery or a power bank (depending on the powerbank, some changes may need to be made). This method requires a USB-B cable. 


The easiest way to power the Arduino Nano 33 BLE Sense board for simple testing is through the MicroUSB connection. The board can be powered by connecting a USB cable to a computer or USB power supply, which is the most common method for powering the device during development and testing.

\section{Power sources Available}

(introduction, compatibility with our board, how long does it last, advantages disadvantages)

Several options are available for powering the Arduino 33 BLE Sense board. To narrow our focus, we will study only the portable ones. This chapter covers three potential sources.
Lithium batterie, Nickel-Metal Hydride (NiMH) batterie, and Power banks.

Choosing the right type is vital as they all have unique advantages and disadvantages.


%The voltage needed for the board depends on the specific port used for power.
%
%1. USB Port: The USB port on the board provides a regulated 5 volts (5V) DC power supply. the USB port provides a regulated 5 volts (5V) DC power supply, and the input voltage supported by the port is between 4.5 volts to 5.5 volts. It is recommended to provide a minimum of 500mA current output for stable operation of the board.
%
%2. Vin Pin: The Vin pin on the board is designed to accept a wide range of input voltages, between 6,2 and 20 volts DC. The voltage on the Vin pin is regulated by the board's onboard buck converter to provide a stable 5V output voltage for powering the board's components.

Now, let's take a look at the three possible power sources:

%VIN pin: The Arduino Nano 33 BLE Sense also has a VIN pin, which can be used to supply power to the device. This pin is connected to the voltage regulator, and it can be used to power the device with an external voltage source.

%(https://www.youtube.com/watch?v=c03UuefFB3A&list=LL&index=12&t=160s&ab_channel=ProgrammingElectronicsAcademy)

\subsection{Nickel-metal hydride (NiMH)}

Nickel-metal hydride (NiMH) batteries are a popular choice because of their rechargeability; the most common is a nominal voltage of 1.2 volts per cell, which means you will have to use several of them.

NiMH batteries offer reliability and cost-effectiveness for various applications including digital cameras, flashlights, remote controls, and handheld gaming consoles thanks to their long cycle life.

NiMH batteries are a versatile choice for numerous users, especially in devices with low power consumption that require extended battery life. However, they may not be the optimum decision for high-power, high-drain devices. They might lack enough energy to keep the item running for extended durations. Moreover, NiMH batteries come with certain drawbacks. As an illustration, their charge diminishes over time even when not being used since they tend to self-discharge. Additionally, they are not as powerful and have lesser capacity compared to other battery kinds, like lithium-ion.\cite{Babu:2023}

Notwithstanding these drawbacks, NiMH batteries remain a top choice for specific applications. As a case in point, they can be charged rapidly which proves useful when the battery is quickly depleted as seen with portable electronics and communications. So, if you require a battery that can be charged rapidly and is appropriate for low-drain devices, NiMH batteries are a suitable choice. However, if your needs include a battery for high-power and high-drain devices, it might be worthwhile to explore better choices. Take lithium-ion batteries as an example, they are more powerful and have a higher capacity.\cite{Babu:2023}

The bigger version of them has been used in "prototypes" of hybrid cars, the Toyota Prius, models II and III came with a fast, heavy NiMH battery with around 200 V and around 6.5 Ah stocked. These batteries achieve enormous lifespans and cycle numbers in such cars. However, for plug-in hybrids and newer hybrid models, Toyota and other manufacturers have switched to lithium to save weight, after apparently solving the problems associated with this technology. \cite{Scherer:2022}

NiMH batteries are considered a safer option for safety-sensitive applications due to the inherent non-flammability of their aqueous electrolyte. Additionally, they also perform better at low temperatures compared to lithium-ion batteries because ions are more mobile in aqueous electrolytes than in organic electrolytes.
\cite{Babu:2023}



\subsection{Lithium batteries}

%General 

Highly preferred for numerous applications owing to their notable features such as extended lifespan, a higher power level, and greater energy intensity are lithium batteries. In other words, this denotes their capacity to store ample amounts of energy in a tiny space while delivering high-powered outputs swiftly and operating for prolonged periods without needing battery replacements \cite{Babu:2023}.

Energy densities up to 141 Wh kg\textsuperscript{-1} and power densities up to 20,600 W kg\textsuperscript{-1} can be achieved from an aqueous lithium-ion battery with a single LiVPO\textsubscript{4}F material as both cathode and anode when using a “water-in-salt” gel polymer electrolyte. This battery also remains flexible during >4000 cycles. This exceeds reported aqueous energy storage devices by far at the same power level \cite{Yang:2017}.


If the cell is not fully discharged and has a residual voltage of around 2.5V, it is charged using a CC/CV approach - first constant current with a charging voltage and ending the charging process when a defined voltage is reached. \cite{Scherer:2022}

%work method
A lithium-ion battery consists of four main components: cathode, anode, electrolyte, and separator. The cathode determines the capacity and voltage of the battery by using lithium oxide as an active material. The anode allows the flow of electric current through the external circuit while enabling reversible absorption/emission of lithium ions released from the cathode. The electrolyte enables the movement of only lithium ions between the cathode and anode while the separator functions as a physical barrier to prevent the direct flow of electrons and carefully allows only the ions to pass through\cite{Samsung:2016,Li:2021}. %(https://www.samsungsdi.com/column/technology/detail/55272.html) 
These components work together to generate electricity through chemical reactions of lithium, which results in the movement of lithium ions and electrons generating electricity. (See \ref{Li-ion battery composition} for a better understanding)

\begin{figure}
    \centering
    \includegraphics[width=9cm]{PowerSource/LithiumBatteryComposition}
    \caption{Li-ion battery composition  \cite{Samsung:2016}}
    \label{Li-ion battery composition}
\end{figure}
%Problems 
Considering the potential impact of moisture on battery safety and performance is crucial. High humidity may cause corrosion and damage to the nail penetration of high energy-density lithium-ion batteries resulting in compromised battery safety.
Their life cycle may deteriorate due to several reasons such as rate flow, current, number and sequence while charging/discharging, temperature variations, and storage methods. The recommended temperature range for Li-ion operation is between 15°C and 35°C [European]. 

Several effects can result from exposing lithium-ion batteries to low temperatures including the loss of both capacity and power as well as life degradation. Furthermore, safety hazards may arise alongside charging difficulties or an unbalanced charge distribution and supplementary expenditure. The study of correlations between these effects is ongoing, along with the exploration of potential solutions to enhance the low-temperature operating scenarios of lithium-ion batteries \cite{Vidal:2019}.

Compared to other types of batteries, Li-ion batteries exhibit lower thermal stability, which has resulted in numerous incidents of battery fires in various applications, such as mobile phones, electric vehicles, and airplanes, in recent years. \cite{Kong:2018,Ouyang:2019,Sun:2020}

To prevent damage to lithium batteries due to overheating during charging, high-quality chargers and battery packs often include thermal monitoring, typically in the form of a thermal resistor. It's safer for us to use high-quality cells because overheating can cause a cell to expand or even become hot enough to ignite. Figure \ref{damaged_lithium_battery} shows damaged lithium batteries that have expanded from their original flat shape to an almost round shape. \cite{Scherer:2022}

%compatibility with our device
To operate the Arduino Nano 33 BLE sense, a minimum input voltage of 4.8V is required. However, most Lithium batteries available in the market only provide a maximum of 4.2V and a minimum of 3.7V. To address this issue, two potential solutions are available.

The first solution is to connect two batteries in series, which would enable us to meet the minimum input requirement and will as a bonus result in longer battery life for the device. Although, it should be noted that this solution would increase the weight, cost, and size of the device.

The second solution entails utilizing a boost converter to attain the appropriate voltage, which is crucial in ensuring that the board is powered effectively. This approach would enable us to achieve the necessary voltage without significant additional weight or volume to the device.


\begin{figure}
    \centering
    \includegraphics[width=0.5\textwidth]{PowerSource/DamagedLithiumBattery}
    \caption{Damaged Lithium battery} %\cite{}
    \label{damaged_lithium_battery}
\end{figure}

\subsection{Powerbanks}

There are colorful types of power banks available in the request, differing in size and power capacity. Some are small and movable, ideal for charging smartphones on- the go, while others are larger and can charge multiple biases contemporaneously or indeed power laptops. The power capacity of power banks also varies, with some offering a few thousand mAh while others can go up to tens of thousands of mAh. It's important to consider your charging requirements and device comity when choosing a power bank.

Power banks are an accessible and movable way to charge electronic bias on the go, they're movable and can be fluently carried around, extending the operating time of mobile bias without the need to charge them through a bowl that's connected to the mains, and they can serve as a provisory power source for electronic bias in situations where there's no access to a power outlet or during power outages. They're also useful for trial and educational purposes \cite{Holz:2020}.  

The size of a power bank is generally commensurable to its capacity, meaning that power banks with advanced capacities tend to be larger in size, utmost power banks use lithium-ion batteries, which can hold a charge for days. Still, the charging miracle that power banks use is time-consuming, which can be a debit for druggies who need to charge their bias snappily \cite{Sunanda:2020}.

Luc Lemmens emphasized in a laboratory trial featured in the Elektor special magazine 2022 that power banks differ in the quantum of energy they can store, and that the energy that a power bank can give should be of interest from a physical point of view. The typical affair voltage and amperage of a marketable power bank can vary depending on the specific model and capacity.  

One issue that may arise when using a power bank is that if the connected device has a too low current draw, the power bank may be unfit to serve duly or will shut itself down. Power banks are designed to switch off after many seconds as soon as the smartphone is completely charged, and this could only be detected from the affair current. However, it would mean that the phone is charged, if the affair current falls below a certain threshold. 

Given our circumstances, it's likely that the power bank will automatically shut off after a brief period if we don't make any custom variations to it. This is because we're exercising a low-power board \cite{Sunanda:2020}. 
To ensure safety when using power banks, it is important to avoid leaving them in places with high temperatures, such as cars parked under the sun. Charging a power bank while it's hot is also not recommended, so be sure to let it cool down before charging. Taking these precautions can help prevent any potential accidents or damage to your power bank.

Power banks are generally safe to use, but users should be aware of safety concerns such as the safety concerns when testing large station battery banks. Users should also use power banks with caution and according to the manufacturer's instructions to avoid overheating, overcharging, and other safety hazards.


\subsection{How To Calculate Battery Run Time}


The theoretical formula for calculating battery life is: 

\medskip

\PYTHON{Time (h) = Capacity (Ah) / Current (A)}

\medskip

where time is in hours, capacity is given in Ah and current is in A. so a 10 (Ah) battery delivering 1(A) should last for 10 hours, if the current of the powered device is 10A then it would last for 1 hour, and if it is 5A would last for 2 hours. 

Most batteries, especially those with circuits, will not work down to 0 volts and will stop working at a set voltage before the battery is fully drained, usually around 80-90\% of the capacity. 

therefore, depending of the battery´s specifications, the calculation will need to times 0.8-0.9:

\medskip

\PYTHON{Hours Time(h) = Capacity(mAh)/Current(mA)*0,9}

\medskip

If we assume that the current consumption of the Arduino board is 20mA and that we are using 4 AA Nickel-Metallhydrid batteries with a capacity of 2600mAh each then: 

\PYTHON{H= 4*2600mAh/50mA *0,9 = 187,2 hrs or 7,8 days (24h)}

To calculate battery life in hours when only Watts are given, it is possible to use the formula: 

\medskip

\PYTHON{Discharging Time=Battery Capacity*Battery Volt/Device Watt.}

\medskip

For example, a 5AH battery with 3V attached to a 10W bulb  5AH*3.7V/10W  would last for 1.85 hours in theory, and 1.66 hours in reality with 90\% power efficiency for Li-ion/LiPo batteries. 






